\section{Appendix}
\label{app:A}

\subsection{ADAS and autonomous vehicle venchmarks}
\label{app:adas}
We provide more details for the different scenarios used for testing ADAS and Autonomous driving control systems.
%

Recall, that each vehicle model in \Simulink\ has several continuous variables including the  $x, y$-coordinates of the vehicle on the road, its velocity, heading, steering angle, etc. 
The vehicle can be controlled by two input signals, namely the throttle (acceleration or brake) and the steering speed.
%
By choosing appropriate values of these input signals, we have defined the following  modes for each vehicle
\begin{inparaenum}[(a)]
	\item \Lmode{cruise}: move forward at constant speed, 
	\Lmode{speedup}: constant acceleration,
	\Lmode{brake}: constant (slow) deceleration,
	\Lmode{em\_brake}: constant (hard).
\end{inparaenum}
We have designed lane switching modes \Lmode{ch\_left} and \Lmode{ch\_right} in which the acceleration and steering are controlled in such a manner that the vehicle switches to its left (resp. right) lane in a certain amount of time. 

For each vehicle, we mainly analyze four variables: absolute position
($sx$) and velocity $vx$ orthogonal to the road direction ($x$-axis),
and absolute position ($sy$) and velocity $vy$ along the road
direction ($x$-axis). The throttle and steering information can be
expressed using the four variables. We will use subscripts to
distinguish between different vehicles.  The following scenarios are
constructed by defining appropriate sets of initial states and
transitions graphs labeled by the modes of two or more vehicles.
%
\begin{description} 
	\item[$\auto{MergeBehind}$:] 
	Initial condition:  Vehicle A is in left and  vehicle B is in the right lane; initial positions and speeds are in some range;  A is in \Lmode{cruise} mode, and B is in \Lmode{cruise} or \Lmode{speedup}.
	%
	Transition graph:  Vehicle A  goes through the mode sequence \Lmode{speedup}, \Lmode{ch\_right}, \Lmode{cruise} with specified intervals of time to transit from mode to another mode. 
	%
	Requirement: A merges behind B within a time bound and maintains at least a given safe separation.
	%%
	%%%%
	%%%%
	%%
	\item[$\auto{MergeAhead}$:] 
	Initial condition: Same as  $\auto{MergeBehind}$ with 
	except that B is in \Lmode{cruise} or \Lmode{brake} mode.
	Transition graph: Same structure as  $\auto{MergeBehind}$ with different  timing parameters.
	%
	Requirement: A merges ahead of B and maintains at least a given safe separation. 
	%%
	%%%%
	%%%%
	%%
	\item[$\auto{AutoPassing}$:]
	Initial condition: Vehicle A behind  B in the same lane, with A in \Lmode{speedup} and B in \Lmode{cruise}; initial positions and speeds are in some range.
	Transition graph:  A goes through the mode sequence \Lmode{ch\_left}, \Lmode{speedup}, \Lmode{brake}, and  \Lmode{ch\_right}, \Lmode{cruise} with specified time intervals in each mode to complete the overtake maneuver. If B switches to
	\Lmode{speedup} before A enters \Lmode{speedup} then
	A aborts and changes back to right lane. If B switches to \Lmode{brake} before A enters \Lmode{ch\_left}, then A should adjust the time to switch to \Lmode{ch\_left} to avoid collision.
	%
	%Car 2 is ahead in right lane and if it is in modes \Lmode{cruise} or \Lmode{brake} the Car 1 proceeds with the maneuver, otherwise if Car 2 accelerates then Car 1 aborts and  goes to  \Lmode{ch\_right} to stay in the right lane.
	Requirement: Vehicle A overtakes B while maintaining minimal safe separation.
	
	\item[$\auto{AEB}$:]
	(Emergency brakes) 
	Initial condition: Vehicle A behind  B in the same lane with A in \Lmode{cruise}, B is stopped (in \Lmode{cruise} mode with velocity $0$). Initial positions and speeds are in some range;
	Transition graph: A transits from \Lmode{cruise} to \Lmode{em\_brake}  over a given interval of time or several disjoint intervals of time.
	Requirement: Vehicle A stops behind B and maintains at least a given safe separation.
	
	\item[$\auto{MergeBetween}$:]
	Initial condition: Vehicle A, B, C are all in the same lane, with A behind B, B behind C, and in the \Lmode{cruise} mode, initial positions and speeds are in some range.
	Transition graph: A goes through the mode sequence \Lmode{ch\_left}, \Lmode{speedup}, \Lmode{brake}, and  \Lmode{ch\_right}, \Lmode{cruise} with specified time intervals in each mode to overtake B. C transits from \Lmode{cruise} to \Lmode{speedup} then transits back to \Lmode{cruise}, so C is always ahead of A.
	Requirement: Vehicle A merges between B and C and any two vehicles maintain at least a given safe separation.
\end{description}



\subsection{Automatic transmission control}
\label{ssec:gear}
We provide some details about the Automatic transmission control benchmark that we have modeled as a hybrid system that combine white-box and black-box components and we have verified using \toolname's safety verification algorithm.

This is a slightly modified version of the Automatic Transmission model provided by \Mathworks\ as a \Simulink\ demo \cite{Matlab_trans}. 
It is a model of an automatic transmission controller that exhibits both continuous and discrete behavior.
The model has  been previously  used by  S-taliro~\cite{S-Taliro} for falsifying certain requirements. We are not aware of any verification results for this system.
% cite all relevant papers by others; Jyo? Sanjit?

For our experiments, we made some minor modifications to  the \Simulink\ model to create the hybrid system  $\auto{ATS}$. 
This allows us to simulate the vehicle from any one of the four modes, namely,  \Lmode{gear1}, \Lmode{gear2}, \Lmode{gear3} and \Lmode{gear4}.
%
Although the system has many variables, we are primarily  interested in the car Speed ($v$), engine RPM (Erpm), impeller torque ($T_i$), output torque ($T_o$), and transmission RPM (Trpm), and therefore, use simulations that record these. 
%
Transition graph of  $\auto{ATS}$ encodes 
transition sequences and intervals for
shifting from \Lmode{gear1} through to \Lmode{gear4}.
Requirement of interest is that the engine RPM is less than a specified maximum value, which in turn is important for limiting the thermal and mechanical stresses on the cylinders and camshafts. Typical unsafe set $\U_t$ could be  Erpm $>4000$.
%
%In the safety verification experiments reported in Section~\ref{sec:reachexp}, we use several different initial sets the initial set  
%$v=0, T_i = 394.05, T_o = 52.92,$ Trpm $= 0$ and 
%Erpm to be within the specified.

%\subsection{Proof for \propref{linear-sep-learn}}
%\label{sec:proof}
%\propref{linear-sep-learn}. Let $\epsilon, \delta \in \plreals$. If
%$k \geq \frac{1}{\epsilon}\ln\frac{1}{\delta}$ then, with probability
%$\geq 1-\delta$, the above algorithm finds $(a,b)$ such that
%$\mathsf{err}_{{\cal D}}(a,b) < \epsilon$.
%
%%
%\begin{proof}
%The result follows from the PAC-learnability of concepts with low
%VC-dimension~\cite{KearnsVazirani}. However, since the proof is very
%simple in this case, we reproduce it here for completeness. Let $k$ be
%as in the statement of the proposition, and suppose the pair $(a,b)$
%identified by the algorithm has error $> \epsilon$. We will bound the
%probability of this happening.
%
%Let $B = \{(x,y)\: |\: x > ay+b\}$. We know that ${\cal D}(B) >
%\epsilon$. The algorithm chose $(a,b)$ only because no element from
%$B$ was sampled in Step~\ref{alg:linsep1}. The probability that this
%happens is $\leq (1-\epsilon)^k$. Observing that $(1-s) \leq e^{-s}$
%for any $s$, we get $(1-\epsilon)^k \leq e^{-\epsilon k} \leq e^{-\ln
%  \frac{1}{\delta}} = \delta$.  This gives us the desired result.
%\end{proof}

\subsection{Safety verification algorithm}
\label{appendix:safetyveri}

The safety verification algorithm is shown in~\ref{alg:safetyveri}. It proceeds along the line of the simulation-based verification algorithms presented in~\cite{DMV:EMSOFT2013, FanMitra:2015df, DuggiralaMV:2015c2e2}.

\begin{algorithm}
	\caption{$\mathit{VerifySafety}(\H,\U)$ verifies safety of hybrid system $\H$ with respect to unsafe set $\U$.}
	\label{alg:safetyveri}
	\SetKwInOut{Input}{input}
	\SetKwInOut{Initially}{initially}
	%\Input {$\H = \langle \L,\Theta, G, \TL\rangle, \U$}
	\Initially{$\I.push(Partition(\Theta))$}
	\While {$\I \neq \emptyset$}
	{	
		$S \gets \I.pop()$\;
		$RS \gets \GraphReach(\H)$ \;
		\uIf {$RS \cap \U = \emptyset$}
		{ continue\;}
		\uElseIf {$\exists (x,l,t) \in RT$ s.t. $\langle RT,v \rangle \in RS$ and $(x,l,t) \subseteq \U $}
		{\Return UNSAFE, $\langle RT,v \rangle$}
		\Else
		{
			{$I.push (Partition(S))$ \;}	
			{Or, $G \gets RefineGraph(G)$ \;}
		}
	}
	\Return SAFE
	%\Return $\R$ \;
\end{algorithm}

%\subsection {Proof for \propref{prop:check-sim}}
%\label{sec:proof2}
%\propref{prop:check-sim}. 
%Given graphs $G_1$ and $G_2$, and mode map $\modemap$, checking if
%there is a forward simulation from $G_1$ to $G_2$ is in polynomial
%time.
%
%\begin{proof}
%The result can be seen to follow from the algorithm for checking timed
%simulations between timed automata~\cite{cerans} and the
%correspondence between one-clock timed automata; the fact that the
%automata have only one clock ensures that the region construction is
%poly-sized as opposed to exponential-sized. However, in the special
%case of transition graphs there is a more direct algorithm which does
%not involve region construction that we describe here.
%
%Observe that if $\{R_i\}_{i\in I}$ is a family of forward simulations
%between $G_1$ and $G_2$ then $\cup_{i \in I} R_i$ is also a forward
%simulation. Thus, like classical simulations, there is a unique
%largest forward simulation between two graphs that is the greatest
%fixpoint of a functional on relations over states of the transition
%graph. Therefore, starting from the relation $\V_1 \times \V_2$, one
%can progressively remove pairs $(v,u)$ such that $v$ is not simulated
%by $u$, until a fixpoint is reached. Moreover, in this case, since
%$G_1$ is a DAG, one can guarantee that the fixpoint will be reached in
%$|\V_1|$ iterations.
%\end{proof}


%\subsection{Graphs and reachtubes for graph simulation}
%\label{appendix:figure_aed}
%
%\begin{figure}[ht]
%    \begin{subfigure}[b]{0.3\textwidth}
%        \includegraphics[width=\textwidth]{Figures/brake_2_modes.png}
%        \caption{\small Transition graph $G_1$.}
%        \label{fig: em_brake_graphsA}
%    \end{subfigure}
%    \begin{subfigure}[b]{0.3\textwidth}
%        \includegraphics[width=\textwidth]{Figures/brake_3_modes.png}
%        \caption{Transition graph $G_2$.}
%        \label{fig: em_brake_graphsB}
%    \end{subfigure}
%    \begin{subfigure}[b]{0.35\textwidth}
%	\includegraphics[width=\textwidth]{Figures/Abstraction_reachTube.eps}
%	\caption{$\auto{AEB}$ Reachtubes.}
%	%\caption{Reachtubes of graphs in Figure \ref{fig: em_brake_graphs}. Red tubes corresponds to Figure \ref{fig: em_brake_graphsA}, blue and gray tubes corresonds to Figure \ref{fig: em_brake_graphsB}}
%	\label{fig:abstraction}
%    \end{subfigure}
%\caption{\small Graphs and reachtubes for the Automatic Emergency Braking $\auto{AEB}$ system.}
%\label{fig: em_brake_graphs}
%\end{figure}
